\documentclass[11pt]{article}
\usepackage{icmlko}
\begin{document}

\title{두 배냐 예순 배냐 --- 라벨이 로그였다는 것을 알고 다시 재니, 실무의 판은 두 배 차이이고 우리는 거기서 $57\%$ 다}
\author{월드모델 실험실}
\date{노트 533 · 2026년 8월 4일}
\maketitle

\begin{abstract}
노트 532 가 ``$0\!\sim\!5$ 점 쌍이 전체의 얼마인가 --- 그만큼이 원리상 못 푸는
몫''을 남겼다. 백분위로 세면 $9.2\%$ 인데 그것은 \emph{정의상} 늘 있다.
사업적으로 뜻이 있는 물음은 \textbf{실제 성과가 몇 배 차이인가}다.

\begin{center}
\begin{tabular}{lrr}
\hline
실제 성과 배수 & 정확도 & 쌍 비율 \\
\hline
$1.0\!\sim\!1.1$ 배 & $51.3\%$ & $4.6\%$ \\
$1.1\!\sim\!1.25$ & $52.4\%$ & $6.3\%$ \\
$1.25\!\sim\!1.5$ & $54.4\%$ & $8.5\%$ \\
$1.5\!\sim\!2$ & $57.4\%$ & $12.5\%$ \\
$2\!\sim\!5$ & $64.5\%$ & $29.1\%$ \\
$5$ 배 이상 & $\mathbf{80.6\%}$ & $39.0\%$ \\
\hline
\end{tabular}
\end{center}

\textbf{후보 다섯 중 $1$ 등과 $2$ 등은 실제로 $\mathbf{1.97}$ 배 차이}이고
그 구간 정확도는 $\mathbf{57.4\%}$ 다. $1$ 등과 $5$ 등은 $\mathbf{61.3}$ 배
차이이고 $\mathbf{80.6\%}$ 다.
\end{abstract}

\begin{figure}[t]
\centering
\includegraphics[width=\textwidth]{figs/f1.pdf}
\caption{\textbf{가} --- 배수별 정확도와 실무의 판. \textbf{나} --- 되돌린
라벨. \textbf{다} --- 읽기.}
\end{figure}

\section{먼저, 내가 단위를 틀렸다}

처음에는 \emph{상대 라벨 차} $r = |y_i - y_j| / \max(y_i, y_j)$ 를 쟀고 사전
등록 셋이 전부 통과했다($r \ge 20\%$ 에서 $80.5\%$, $r < 5\%$ 가 $18.7\%$).

\textbf{그런데 라벨이 로그였다.} 열 도메인이 $\log_{10}$(관심 수 $\cdot$
후원자 수 $\cdot$ 방문자 수 $\cdots$)이고 게임 하나만 자연로그다(노트 210 이
누적 리뷰를 $30$ 일 창으로 바꾸며 \texttt{log1p}). 로그 값에 ``상대 차''를
씌우면 아무것도 뜻하지 않는다.

\begin{quote}
\textbf{검산} --- 되돌린 값이 정수에 얼마나 붙나 보니 아홉 도메인에서
$\mathbf{100\%}$ 다(팝업 $26\%$ $\cdot$ 시장팝업 $44\%$ 는 라벨이 \emph{일평균}
이라 정수가 아닌 것이 맞다). 변환 추정이 정확하다.
\end{quote}

되돌린 중앙값은 이렇다 --- 아이돌 초동 $206{,}700$ $\cdot$ 세계애니 한줄평
$38{,}380$ $\cdot$ 웹툰 관심 $32{,}100$ $\cdot$ 도서 판매지수 $14{,}790$
$\cdot$ 게임 리뷰 $3{,}872$ $\cdot$ 시장팝업 일방문 $1{,}310$ $\cdot$ 만화
서재 $1{,}133$ $\cdot$ 팝업 일방문 $458$ $\cdot$ 모바일 평가 $169$ $\cdot$
애니 한줄평 $63$ $\cdot$ 펀딩 후원자 $37$.

\subsection*{되돌리니 겸손해진다}

\begin{center}
\begin{tabular}{lrr}
\hline
 & 틀린 단위 & \textbf{맞는 단위} \\
\hline
``뜻이 있는 차이''에서의 정확도 & $80.5\%$ & $\mathbf{73.1\%}$ ($2$ 배 이상) \\
``거의 같은'' 쌍의 비율 & $18.7\%$ & $\mathbf{4.6\%}$ ($1.1$ 배 미만) \\
\hline
\end{tabular}
\end{center}

사전 등록 ②는 맞는 단위에서도 통과하지만($73.1 \ge 72$) 여유가 사라졌고,
③은 \textbf{반증}이다 --- ``거의 같은'' 쌍은 $15\%$ 가 아니라 $4.6\%$ 다.

\section{덮음과 함께 (노트 532 조항)}

\begin{center}
\begin{tabular}{lrr}
\hline
문턱 & 덮음 & 정확도 \\
\hline
$1.10$ 배 이상 & $95.4\%$ & $67.9\%$ \\
$1.25$ 배 이상 & $89.1\%$ & $69.0\%$ \\
$1.50$ 배 이상 & $80.6\%$ & $70.6\%$ \\
$\mathbf{2.00}$ \textbf{배 이상} & $\mathbf{68.1\%}$ & $\mathbf{73.1\%}$ \\
\hline
\end{tabular}
\end{center}

\section{이어 붙이면 --- 세 번째로 같은 답}

\begin{center}
\begin{tabular}{lrr}
\hline
 & 실제 배수 & 그 구간 정확도 \\
\hline
후보 다섯 중 $1$ 등 대 $2$ 등 & $\mathbf{1.97}$ 배 & $\mathbf{57.4\%}$ \\
후보 다섯 중 $1$ 등 대 $5$ 등 & $61.3$ 배 & $80.6\%$ \\
\hline
\end{tabular}
\end{center}

\begin{quote}
\textbf{모델의 실력은 성과 배수를 따라간다.} 실무의 판($1$ 등 대 $2$ 등)은
\emph{두 배} 차이라 어렵고, 거르는 판($1$ 등 대 $5$ 등)은 \emph{예순 배} 차이라
쉽다. 노트 528 이 결정 단위에서 ``나쁜 것 거르기 $3.22$ 배 대 $1$ 등 집기
$2.00$ 배''로, 노트 $531 \cdot 532$ 가 ``접전에서 신호가 침묵''으로 말한 것과
\textbf{같은 기계다. 세 번째로 같은 답이 나왔다.}
\end{quote}

\textbf{도메인 차이도 크다} --- $1$ 등 대 $5$ 등 배수가 모바일 $296$ 배 $\cdot$
게임 $117$ 배인데 \textbf{만화 $3.5$ 배 $\cdot$ 세계애니 $3.7$ 배}다. 그 둘은
성과 분산 자체가 작아서 어떤 모델도 가르기 어렵다.

\section*{조항}

\begin{quote}
\textbf{지표를 만들기 전에 라벨이 어떤 변환 위에 있는지 확인한다.} 그리고
결과는 \textbf{되돌린 단위}로 적는다 --- ``상대차 $20\%$''는 아무도 못 읽지만
``방문자 두 배''는 누구나 읽는다. 되돌린 값을 \emph{정수 근접}으로 검산하면
변환 추정이 맞는지 공짜로 확인된다.
\end{quote}

\section*{남은 것}

\begin{itemize}
\item \textbf{만화 $\cdot$ 세계애니는 성과 분산이 작다}($1$--$5$ 등 $3.5$ 배).
이 둘의 낮은 정확도는 모델 탓이 아닐 수 있다 --- 노트 505 의 자기 천장과
겹쳐 봐야 한다.
\item \textbf{두 배 문턱이 사업적으로 맞나} --- 방문자 두 배는 큰 차이지만
후원자 $37$ 명 대 $74$ 명은 둘 다 실패다. \emph{절대} 수준을 같이 봐야 할 수
있다.
\item \textbf{덱을 고쳐야 한다} --- 지금 ``$67.4\%$''만 적혀 있고 배수 축이 없다.
\end{itemize}

\end{document}
